\section{Introduction}
\label{submission:sec:intro}

The exponential growth of specialized deep learning models has fueled intense interest in weight-space consolidation, commonly known as \emph{model merging} or \emph{weight ensembling} \cite{wortsman2022model, ilharco2022editing}. By directly combining the parameters of multiple task-specific expert neural networks fine-tuned from a shared pre-trained backbone, model merging enables multi-task capabilities without the prohibitive cost of joint training or the increased inference latency associated with multi-branch routing architectures. Techniques such as Task Arithmetic \cite{ilharco2022editing}, Git Re-Basin \cite{ainsworth2022git}, ZipIt! \cite{stoica2023zipit}, and AdaMerging \cite{yang2023adamerging} have demonstrated remarkable success in retaining expert capabilities within a single consolidated parameter set.

However, deployment on real-world edge hardware imposes strict constraints on both memory footprint and computational latency, necessitating the use of post-training quantization (PTQ) to compress weights to 8-bit or 4-bit integer representations \cite{nagel2021whitepaper}. When merged models are quantized naively post-hoc---a process we term \emph{Merge-then-Quantize (M-then-Q)}---the non-linear discretization noise introduced by rounding functions often triggers a catastrophic degradation of multi-task performance, particularly at low bit-widths (e.g., 4-bit) \cite{qmerge2025}. To address this, the emerging paradigm of \emph{Quantization-Aware Model Merging} (exemplified by Q-Merge \cite{qmerge2025}) has proposed optimizing layer-wise merging coefficients directly under a specific quantization operator using Straight-Through Estimators (STE) \cite{bengio2013ste}. Guided by unsupervised objectives like prediction entropy minimization \cite{wang2020tent}, these methods claim to find "near-lossless" low-bit merging configurations using tiny calibration streams.

In this work, we approach these "state-of-the-art" claims with methodological skepticism. In deep learning, overly optimistic claims often arise from narrow, idealized evaluation setups that obscure critical confounding variables and fragile assumptions. We identify several unstudied assumptions in the current evaluation protocols of quantization-aware merging:
\begin{enumerate}
    \item \textbf{Quantization-Operator Monomorphism:} Prior evaluations assume that merging coefficients optimized under a simulated "fake" quantization operator (e.g., Symmetric Per-Channel) will be deployed on a hardware platform with an identical mathematical representation.
    \item \textbf{Calibration Stream Purity:} Current frameworks assume the availability of highly pristine, perfectly class-balanced, and static test-time calibration data streams, which rarely reflect the messy, corrupted, and skewed realities of streaming edge inference.
    \item \textbf{STE Gradient Path Fidelity:} The straight-through gradient approximation is treated as a high-fidelity guide for weight-space search, ignoring its potential to overfit to discretization noise or get trapped in brittle local minima.
\end{enumerate}

To expose these blind spots, we present a rigorous, independent \emph{Multi-Axial Robustness Audit and Methodological Deconstruction} of Q-Merge. Using a standardized pre-trained Vision Transformer backbone (\texttt{timm ViT-Tiny}) \cite{dosovitskiy2020vit} with four distinct classification heads trained on MNIST, FashionMNIST, CIFAR-10, and SVHN, we systematically dissect the behavior of quantization-aware model merging along four distinct axes:
\begin{itemize}
    \item \textbf{Axis 1: Calibration Stream Size Sweep.} We sweep the calibration size $N \in \{1, 4, 16, 64\}$ per task to determine where transductive overfitting to discretization noise occurs and whether larger sample sizes resolve performance collapse.
    \item \textbf{Axis 2: Cross-Schema Generalization Matrix.} We introduce a cross-schema evaluation framework where coefficients are optimized under a source quantization schema $Q_{\text{opt}}$ but deployed under five different hardware-relevant target schemas $Q_{\text{eval}}$ (e.g., symmetric/asymmetric, per-tensor/per-channel, and double quantization).
    \item \textbf{Axis 3: Spatial Regularization \& Derivative-Free Search.} We evaluate whether spatial smoothing regularizers (like Total Variation) or alternative derivative-free black-box optimization (1+1 Evolution Strategy) can mitigate overfitting to the simulated operator's rounding boundaries.
    \item \textbf{Axis 4: Stream Distortion and Skew Robustness.} We stress-test the optimization under non-idealized streaming conditions, specifically injecting out-of-distribution (OOD) Gaussian corruptions and severe class imbalance (Gini skew).
\end{itemize}

Our audit yields several highly surprising and methodologically significant findings that challenge current trends:
\begin{itemize}
    \item \textbf{Catastrophic Cross-Operator Overfitting:} Merging coefficients overfit intensely to the exact mathematical operator used during optimization. Moving from a channel-wise source to a tensor-wise target causes accuracy to collapse back to random-guess levels (~10\%), revealing that the optimized weights represent a highly brittle local configuration rather than a genuine weight-space alignment.
    \item \textbf{The Superiority of Full-Precision Search (Quantized AdaMerging):} We discover that unquantized optimization in FP16 followed by post-hoc quantization consistently and substantially outperforms Q-Merge's direct quantization-aware optimization via STE ($30.00\%$ vs $26.25\%$). This reveals that direct low-bit optimization under quantization constraints introduces substantial gradient noise (via STE approximations) that damages weight-space search, whereas full-precision optimization is more robust.
    \item \textbf{The Failure of unconstrained STE:} Unconstrained STE optimization fails to consistently beat a naive M-then-Q baseline across varying sample sizes, and empirical sweeps prove that this instability under hard rounding thresholds cannot be resolved by simple learning rate tuning.
    \item \textbf{Stochastic Search vs. Biased Gradients:} A derivative-free 1+1 Evolution Strategy (1+1 ES) outperforms gradient-based STE on the source schema ($20.75\%$ vs $17.88\%$), highlighting the limits of relying on artificial straight-through gradients in highly non-convex quantized landscapes. However, because 1+1 ES is a powerful black-box searcher, it finds a highly customized and fragile configuration on the source operator's rounding boundaries, leading to even more severe cross-schema overfitting when evaluated on a mismatched target.
    \item \textbf{Label Skew Vulnerability:} Standard unsupervised entropy minimization is blind to label distribution. Under severe class skew, the optimized weights collapse, completely destroying the decision boundaries of underrepresented classes.
\end{itemize}

By highlighting these limitations, we do not seek to dismiss quantization-aware model merging, but rather to establish a more honest, rigorous, and deployment-realistic evaluation standard. The rest of this paper is structured as follows: Section~\ref{submission:sec:related} contextually situates our work; Section~\ref{submission:sec:method} formalizes the mathematics of model merging, quantization, and our cross-operator audit; Section~\ref{submission:sec:experiments} presents our extensive empirical results and multi-axial analysis; and Section~\ref{submission:sec:conclusion} concludes with critical recommendations for future research.
